\section{Техническое задание}
\subsection{Основание для разработки}

Основанием для разработки является задание на выпускную квалификационную работу бакалавра на тему "<Программно-информационная система обнаружения дефектов дорожного покрытия">.

\subsection{Цель и назначение разработки}

Целью разработки является создание программно-информационной системы, предназначенной для автоматизированного обнаружения дефектов дорожного покрытия на основе анализа цифровых изображений с использованием методов компьютерного зрения и нейросетевых алгоритмов.

Назначением системы является повышение скорости выявления дефектов дорожного полотна, а также снижение трудозатрат при визуальном контроле состояния автомобильных дорог.

Основными задачами разработки являются:

\begin{itemize}
	\item анализ существующих методов обнаружения дефектов дорожного покрытия;
	\item выбор алгоритма детекции объектов;
	\item разработка программной архитектуры системы;
	\item реализация модуля обработки изображений;
	\item обучение и интеграция модели нейросети;
	\item разработка графического пользовательского интерфейса;
	\item реализация средств визуализации результатов;
	\item тестирование системы.
\end{itemize}

\subsection{Требования к функциональности}

\subsubsection{Загрузка изображений}

Система должна обеспечивать:

\begin{itemize}
	\item выбор одного или нескольких файлов изображений;
	\item предварительную проверку поддерживаемого формата;
	\item отображение миниатюр загруженных изображений;
	\item отображение информации о количестве загруженных изображений.
\end{itemize}

Загрузка должна осуществляться посредством стандартного диалогового окна операционной системы.

\subsubsection{Настройка параметров анализа}

Система должна предоставлять пользователю возможность настройки параметров обнаружения дефектов, включая:

\begin{itemize}
	\item изменение порога достоверности обнаружения;
	\item выбор директории временного хранения;
	\item выбор режима сортировки обработанных изображений.
\end{itemize}

\subsubsection{Обработка изображений}

Подсистема анализа должна обеспечивать:

\begin{itemize}
	\item применение обученной нейросетевой модели;
	\item последовательную обработку загруженных изображений;
	\item обнаружение объектов, относящихся к дефектам покрытия;
	\item формирование координат ограничивающих рамок;
	\item вычисление вероятности принадлежности объекта к классу.
\end{itemize}

Во время обработки пользователь должен получать информацию о текущем прогрессе выполнения.

\subsubsection{Визуализация результатов}

Система должна обеспечивать:

\begin{itemize}
	\item отображение обработанных изображений;
	\item визуальное выделение найденных дефектов;
	\item вывод статистической информации;
	\item отображение уровня достоверности обнаружения.
\end{itemize}

\subsubsection{Сохранение и восстановление результатов}

Система должна поддерживать:

\begin{itemize}
	\item сохранение обработанных изображений;
	\item формирование текстового отчёта;
	\item загрузку сохранённых результатов;
	\item возможность сортировки вновь открытых результатов.
\end{itemize}

Также должно быть возможно:
\begin{itemize}
	\item очистить временные данные результатов анализа;
	\item выбрать каталог для хранения временных результатов.
\end{itemize}

\subsubsection{Настройка пользовательского интерфейса}

Система должна предоставлять:

\begin{itemize}
	\item переключение между светлой и тёмной темами;
	\item автоматическое обновление визуального оформления при изменении размера окна;
	\item сохранение выбранной темы между сеансами работы.
\end{itemize}

\subsection{Требования к интерфейсу}

Интерфейс программно-информационной системы должен представлять собой графическое настольное приложение с интуитивно понятным пользовательским управлением.

Основные элементы интерфейса должны включать:

\begin{itemize}
	\item главное окно приложения;
	\item область отображения загруженных изображений;
	\item область отображения результатов анализа;
	\item меню управления обработкой;
	\item меню выбора файлов, сохранения и загрузки результатов;
	\item меню изменения вида главного окна и порядка выведенных результатов;
	\item панель статистики.
\end{itemize}

\subsection{Моделирование вариантов использования}

В системе выделяется одно действующее лицо: оператор системы --- выполняет загрузку изображений, запуск анализа и просмотр результатов обработки.

Перечень основных вариантов использования:

\begin{enumerate}
	\item Загрузка изображений.
	\item Настройка параметров обнаружения.
	\item Запуск анализа изображений.
	\item Просмотр результатов обнаружения.
	\item Сортировка результатов.
	\item Сохранение результатов.
	\item Формирование отчёта.
	\item Загрузка ранее сохранённых результатов.
	\item Очистка временных данных.
	\item Переключение темы интерфейса.
\end{enumerate}

Диаграмма прецедентов представлена на рисунке \ref{DiagramVariant:image}.

\begin{figure}[H]
	\center{\includegraphics[width=0.8\linewidth]{DiagramVariant}}
	\caption{Диаграмма прецедентов для пользователя}
	\label{DiagramVariant:image}
\end{figure}

\subsubsection{Вариант использования «Загрузка изображений»}

\textbf{Заинтересованные лица и их требования:} оператор желает загрузить одно или несколько изображений дорожного покрытия для последующего анализа.

\textbf{Предусловие:} приложение запущено, отображается зона вывода входных изображений.

\textbf{Постусловие:} выбранные изображения отображаются в интерфейсе системы.

\textbf{Основной успешный сценарий:}

\begin{enumerate}
	\item Оператор нажимает кнопку загрузки изображений.
	\item Система открывает диалог выбора файлов.
	\item Оператор выбирает один или несколько файлов.
	\item Изображения по ходу загрузки отображаются в зоне вывода входных изображений.
\end{enumerate}

\subsubsection{Вариант использования «Настройка параметров обнаружения»}

\textbf{Заинтересованные лица и их требования:} оператор желает настроить параметры анализа для повышения точности обнаружения.

\textbf{Предусловие:} изображения загружены.

\textbf{Постусловие:} новые параметры сохранены и готовы к использованию.

\textbf{Основной успешный сценарий:}

\begin{enumerate}
	\item Оператор открывает меню обработки.
	\item Во вкладке выбора порога обнаружения изменяет значение.
	\item Система сохраняет настройки.
\end{enumerate}

\subsubsection{Вариант использования «Запуск анализа изображений»}

\textbf{Заинтересованные лица и их требования:} оператор желает автоматически обнаружить дефекты дорожного покрытия.

\textbf{Предусловие:} изображения загружены.

\textbf{Постусловие:} изображения обработаны, результаты сохранены во временной папке.

\textbf{Основной успешный сценарий:}

\begin{enumerate}
	\item Оператор открывает меню обработки и нажимает кнопку запуска анализа.
	\item Выполняется обработка изображений, прогресс выводится в строке статистики.
	\item Результаты отображаются в интерфейсе.
\end{enumerate}

\subsubsection{Вариант использования «Просмотр результатов обнаружения»}

\textbf{Заинтересованные лица и их требования:} оператор желает визуально оценить результаты анализа определенного изображения.

\textbf{Предусловие:} результаты обработки выведены на экран.

\textbf{Постусловие:} оператор получает доступ к результатам обработки.

\textbf{Основной успешный сценарий:}

\begin{enumerate}
	\item Оператор выбирает обработанное изображение.
	\item Система отображает изображение с выделенными дефектами.
	\item Отображается количество найденных объектов и вероятности их обнаружения.
\end{enumerate}

\subsubsection{Вариант использования «Сортировка результатов»}

\textbf{Заинтересованные лица и их требования:} оператор желает упорядочить результаты для удобства анализа.

\textbf{Предусловие:} результаты обработки выведены на экран.

\textbf{Постусловие:} результаты отсортированы по выбранному критерию.

\textbf{Основной успешный сценарий:}

\begin{enumerate}
	\item Оператор в меню вида выбирает критерий сортировки.
	\item Система выполняет перестроение списка.
	\item Отображается обновлённый порядок результатов.
\end{enumerate}

\subsubsection{Вариант использования «Сохранение результатов»}

\textbf{Заинтересованные лица и их требования:} оператор желает сохранить результаты анализа.

\textbf{Предусловие:} анализ завершён.

\textbf{Постусловие:} результаты записаны в выбранный каталог.

\textbf{Основной успешный сценарий:}

\begin{enumerate}
	\item Оператор открывает меню и выбирает команду сохранения.
	\item В открывшемся диалоговом окне указывает каталог.
	\item Система сохраняет обработанные изображения и сопутствующие данные.
\end{enumerate}

\subsubsection{Вариант использования «Формирование отчёта»}

\textbf{Заинтересованные лица и их требования:} оператор желает получить текстовый отчёт по результатам анализа.

\textbf{Предусловие:} анализ завершён.

\textbf{Постусловие:} сформирован текстовый файл отчёта.

\textbf{Основной успешный сценарий:}

\begin{enumerate}
	\item Оператор открывает меню и выбирает команду сохранения.
	\item В открывшемся диалоговом окне указывает каталог.
	\item Система сохраняет обработанные изображения и формирует текстовый отчёт.
\end{enumerate}

\subsubsection{Вариант использования «Загрузка сохранённых результатов»}

\textbf{Заинтересованные лица и их требования:} оператор желает повторно открыть ранее сохранённые результаты.

\textbf{Предусловие:} существуют сохранённые данные, файл отчета не поврежден и не изменен.

\textbf{Постусловие:} результаты отображаются в интерфейсе.

\textbf{Основной успешный сценарий:}

\begin{enumerate}
	\item Оператор в меню выбирает открытие обработанных данных.
	\item В открывшемся диалоговом окне указывает каталог.
	\item Система считывает данные.
	\item Результаты отображаются.
\end{enumerate}

\subsubsection{Вариант использования «Очистка временных данных»}

\textbf{Заинтересованные лица и их требования:} оператор желает удалить временные файлы.

\textbf{Предусловие:} во временном каталоге имеются данные.

\textbf{Постусловие:} временные файлы удалены.

\textbf{Основной успешный сценарий:}

\begin{enumerate}
	\item Оператор в меню обработки выбирает очистку временных данных.
	\item Система запрашивает подтверждение.
	\item После подтверждения выполняется удаление.
\end{enumerate}

\subsubsection{Вариант использования «Переключение темы интерфейса»}

\textbf{Заинтересованные лица и их требования:} оператор желает изменить внешний вид интерфейса.

\textbf{Предусловие:} приложение активно.

\textbf{Постусловие:} применяется новая цветовая тема.

\textbf{Основной успешный сценарий:}

\begin{enumerate}
	\item Оператор открывает меню вида.
	\item Выбирает тему оформления.
	\item Система применяет изменения.
\end{enumerate}

\subsection{Требования к надёжности}

Разрабатываемая программно-информационная система должна обеспечивать стабильную работу при обработке изображений, корректно реагировать на ошибки ввода и предотвращать потерю временных результатов.

\textbf{Устойчивость обработки изображений.} Система должна корректно обрабатывать ошибки, возникающие при загрузке файлов неподдерживаемого формата, повреждённых изображений.

\textbf{Контроль выполнения анализа.} При запуске процедуры обнаружения дефектов система должна обеспечивать завершение обработки каждого изображения независимо от ошибок, возникших при анализе других файлов.

\textbf{Сохранность результатов.} Результаты обработки должны сохраняться во временном каталоге до завершения пользовательской сессии либо до их явного удаления пользователем.

\subsection{Требования к программной и технической совместимости}

Система должна корректно функционировать на персональных компьютерах, удовлетворяющих минимальным требованиям к программному обеспечению и вычислительным ресурсам.

\textbf{Операционная система.} Приложение должно обеспечивать корректную работу в среде операционных систем семейства Windows.

\textbf{Используемые библиотеки.} Для функционирования системы требуется наличие следующих программных библиотек:

\begin{itemize}
	\item PySide6 --- для реализации графического пользовательского интерфейса;
	\item OpenCV --- для предварительной обработки изображений;
	\item NumPy --- для выполнения численных операций;
	\item Ultralytics YOLO --- для выполнения детекции дефектов;
\end{itemize}

\textbf{Аппаратные требования.} Минимальные технические характеристики вычислительного устройства:

\begin{itemize}
	\item процессор с тактовой частотой не менее 2 ГГц;
	\item оперативная память не менее 8 ГБ;
	\item свободное место на диске не менее 4 ГБ;
	\item графический адаптер (видеокарта).
\end{itemize}

\textbf{Форматы входных данных.} Система должна поддерживать обработку изображений форматов JPG, JPEG и PNG.

\textbf{Совместимость модели.} Для корректной работы требуется наличие файла обученной модели нейронной сети формата \texttt{.pt}.

\subsection{Требования к производительности}

Система должна обеспечивать приемлемое время отклика пользовательского интерфейса и достаточную скорость обработки изображений.

\textbf{Время загрузки приложения} не должно превышать 10 секунд.

\textbf{Время обработки одного изображения} размером до 1920x1080 пикселей не должно превышать 5 секунд при использовании рекомендуемой аппаратной конфигурации.

\textbf{Отзывчивость интерфейса.} Во время выполнения обработки система должна отображать индикатор выполнения и предотвращать выполнение конфликтующих операций.

\subsection{Требования к оформлению документации}

Разработка программной документации и программного изделия должна производиться согласно ГОСТ 19.102-77 и ГОСТ 34.601-90. Единая система программной документации.
